\chapter{Triển khai các thành phần}
\label{ch:deployment}

Xây dựng hệ thống theo \textit{Infrastructure as Code (IaC)}, giúp hạ tầng có khả năng tái tạo tự động, dễ mở rộng, dễ triển khai nhiều lần, loại bỏ các sai sót từ thao tác cấu hình thủ công.

\section{Khởi tạo hạ tầng AWS bằng Terraform}
Hạ tầng mạng, cân bằng tải và máy chủ được dựng bằng mã nguồn Terraform trong thư mục \texttt{terraform/}. Trạng thái hệ thống (\textit{Terraform State}) được lưu trữ tập trung và dùng khóa trạng thái (\textit{State Locking}) trên AWS S3 Bucket và DynamoDB:

\begin{lstlisting}[language=bash, caption={Quy trình thực thi khởi tạo hạ tầng Terraform}]
# Khoi tao cac Provider va ket noi S3 Backend
terraform init

# Kiem tra ke hoach thay doi tai nguyen
terraform plan

# Ap dung trien khai ha tang (AWS NLB, 2 Gateways, 3 K8s Nodes, 1 Wazuh Server)
terraform apply -auto-approve
\end{lstlisting}

\subsection{Khởi tạo AWS Network Load Balancer (NLB L4)}
Hệ thống sử dụng AWS NLB để phân phối lưu lượng Internet vào 2 Gateway tại 2 AZ:
\begin{lstlisting}[language=bash, caption={Đoạn mã khai báo AWS NLB trong Terraform}]
resource "aws_lb" "nlb" {
  name               = "ncs-project-nlb"
  internal           = false
  load_balancer_type = "network"
  subnets            = [aws_subnet.public_wan.id, aws_subnet.public_wan_b.id]
}

# 4 Target Groups cho cac cong dich vu: 80, 443, 7443 va 8443
resource "aws_lb_target_group" "tcp" {
  for_each    = { "80" = 80, "443" = 443, "7443" = 7443, "8443" = 8443 }
  port        = each.value
  protocol    = "TCP"
  vpc_id      = aws_vpc.main.id
  target_type = "ip"

  health_check {
    protocol            = "TCP"
    port                = each.value
    interval            = 10
    healthy_threshold   = 2
    unhealthy_threshold = 2
  }
}
\end{lstlisting}

\subsection{Cấu hình 2 card mạng cho máy chủ Gateway}
Mỗi Gateway được trang bị 2 Card mạng ảo (\textit{Elastic Network Interface - ENI}): Card WAN kết nối trực tiếp Internet và Card LAN kết nối vào Subnet nội bộ:
\begin{lstlisting}[language=bash, caption={Đoạn mã cấu hình Dual-ENI trong Terraform}]
resource "aws_instance" "pfsense" {
  ami           = data.aws_ami.ubuntu.id
  instance_type = "t3.small"

  network_interface {
    network_interface_id = aws_network_interface.pfsense_wan.id # Card ens5 (Public)
    device_index         = 0
  }
  network_interface {
    network_interface_id = aws_network_interface.pfsense_lan.id # Card ens6 (Private)
    device_index         = 1
  }
}
\end{lstlisting}

\section{Cấu hình Gateway và hệ thống Suricata IDS/IPS}
Được tự động hóa bằng Ansible playbook \texttt{ansible/site.yml} áp dụng role \texttt{gateway} trên cả hai Gateway:

\subsection{1. Cấu hình định tuyến và IPTables NAT (Chống Asymmetric Routing)}
Kích hoạt chuyển tiếp gói tin Linux, thiết lập DNAT chuyển tiếp cổng và giải quyết triệt để lỗi định tuyến bất đối xứng bằng LAN Masquerade:
\begin{lstlisting}[language=bash, caption={Cấu hình IPTables Port Forwarding và LAN Masquerade trên Gateway}]
# 1. Kich hoat IP Forwarding
sudo sysctl -w net.ipv4.ip_forward=1

# 2. PREROUTING DNAT: Chuyen tiep cong toi dung dich vu noi bo
sudo iptables -t nat -A PREROUTING -i ens5 -p tcp --dport 80 -j DNAT --to-destination 10.0.10.10:80
sudo iptables -t nat -A PREROUTING -i ens5 -p tcp --dport 443 -j DNAT --to-destination 10.0.10.10:443
sudo iptables -t nat -A PREROUTING -i ens5 -p tcp --dport 7443 -j DNAT --to-destination 10.0.10.10:30443
sudo iptables -t nat -A PREROUTING -i ens5 -p tcp --dport 8443 -j DNAT --to-destination 10.0.20.10:55000

# 3. LAN Masquerade: Rewrite Source IP ve IP LAN Gateway (Chong Asymmetric Routing)
sudo iptables -t nat -A POSTROUTING -o ens6 -j MASQUERADE

# 4. WAN Masquerade: Source NAT cho luong Egress cua K8s & Wazuh ra Internet
sudo iptables -t nat -A POSTROUTING -o ens5 -j MASQUERADE

# Luu cau hinh vinh vien tranh bi mat khi reboot
sudo netfilter-persistent save
\end{lstlisting}

\subsection{2. Cấu hình tập rule phát hiện tấn công cho Suricata}
Suricata được cấu hình lắng nghe trực tiếp trên card WAN \texttt{ens5} ở chế độ \texttt{af-packet}. Tập rule tại \texttt{/etc/suricata/rules/local.rules} nhận diện tấn công ứng dụng Web:
\begin{lstlisting}[language=bash, caption={Tập luật phát hiện tấn công mạng trong local.rules của Suricata}]
# Nhan dien SQL Injection (UNION SELECT, OR 1=1)
alert http any any -> any any (msg:"ATTACK: SQL Injection Union Select Detected"; flow:to_server,established; content:"UNION"; nocase; content:"SELECT"; nocase; sid:1000001; rev:1;)
alert http any any -> any any (msg:"ATTACK: SQL Injection OR 1=1 Detected"; flow:to_server,established; content:"or"; nocase; content:"1=1"; sid:1000002; rev:1;)

# Nhan dien Directory Traversal / LFI (/etc/passwd, ../..)
alert http any any -> any any (msg:"ATTACK: Directory Traversal /etc/passwd Detected"; flow:to_server,established; content:"/etc/passwd"; sid:1000003; rev:1;)
alert http any any -> any any (msg:"ATTACK: Directory Traversal Path Manipulation"; flow:to_server,established; content:"../.."; sid:1000004; rev:1;)

# Nhan dien Cross-Site Scripting (XSS)
alert http any any -> any any (msg:"ATTACK: XSS Payload Detected"; flow:to_server,established; content:"<script>"; nocase; sid:1000005; rev:1;)

# Nhan dien do quet tep tin he thong nhay cam
alert http any any -> any any (msg:"ATTACK: Sensitive File Probe (.env/.git)"; flow:to_server,established; content:".env"; sid:1000006; rev:1;)
\end{lstlisting}

\section{Khởi tạo cụm Upstream Kubernetes v1.31 Multi-AZ}
Cụm được triển khai bằng \texttt{kubeadm} phiên bản \textbf{Kubernetes v1.31.14}:
\begin{enumerate}
    \item \textbf{Khởi tạo Control Plane trên Master node}:
    \begin{lstlisting}[language=bash]
sudo kubeadm init --pod-network-cidr=10.244.0.0/16 --apiserver-advertise-address=10.0.10.10
    \end{lstlisting}
    \item \textbf{Cài đặt mạng Flannel CNI}:
    \begin{lstlisting}[language=bash]
kubectl apply -f https://raw.githubusercontent.com/flannel-io/flannel/master/Documentation/kube-flannel.yml
    \end{lstlisting}
    \item \textbf{Join các Worker node}:
    Chạy lệnh \texttt{kubeadm join} trên Worker 1 (\texttt{10.0.10.139} tại AZ-a) và Worker 2 (\texttt{10.0.11.21} tại AZ-b) để hoàn tất liên kết cụm Multi-AZ 3 node.
\end{enumerate}

\section{Triển khai Rancher Management Server}
Rancher Server được triển khai trên cụm K8s thông qua Helm chart của SUSE:
\begin{lstlisting}[language=bash, caption={Lệnh triển khai Rancher Server bằng Helm}]
helm repo add rancher-prime https://charts.rancher.com/server-charts/prime
helm repo update

# Cai dat Rancher vao namespace cattle-system
helm install rancher rancher-prime/rancher \
  --namespace cattle-system \
  --create-namespace \
  --set hostname=rancher.lab \
  --set bootstrapPassword=AdminNCSProject2026! \
  --set replicas=1
\end{lstlisting}

Để tách Rancher khỏi cổng Ingress của ứng dụng thương mại, sử dụng Service \texttt{NodePort} để gắn cổng nội bộ \textbf{\texttt{30443}} với Pod của Rancher:
\begin{lstlisting}[language=yaml, caption={Khai báo NodePort Service cho Rancher Server}]
apiVersion: v1
kind: Service
metadata:
  name: rancher-nodeport
  namespace: cattle-system
spec:
  type: NodePort
  selector:
    app: rancher
  ports:
    - name: https-nodeport
      port: 443
      targetPort: 443
      nodePort: 30443
      protocol: TCP
\end{lstlisting}

\section{Triển khai Ứng dụng Microservices \& Zero-Trust Network Policy}
Toàn bộ 12 dịch vụ của Google Online Boutique được triển khai đồng loạt trong namespace \texttt{online-boutique}.

\begin{lstlisting}[language=yaml, caption={Chính sách NetworkPolicy Zero-Trust bảo vệ PaymentService}]
apiVersion: networking.k8s.io/v1
kind: NetworkPolicy
metadata:
  name: paymentservice-zero-trust
  namespace: online-boutique
spec:
  podSelector:
    matchLabels:
      app: paymentservice
  policyTypes:
    - Ingress
    - Egress
  ingress:
    # Chi chap nhan ket noi den tu CheckoutService qua cong gRPC 50051
    - from:
        - podSelector:
            matchLabels:
              app: checkoutservice
      ports:
        - protocol: TCP
          port: 50051
  egress:
    # Chi cho phep Egress giai ma DNS CoreDNS noi bo
    - to:
        - namespaceSelector: {}
          podSelector:
            matchLabels:
              k8s-app: kube-dns
      ports:
        - protocol: UDP
          port: 53
\end{lstlisting}

\section{Cài đặt Wazuh SIEM trên Máy chủ Chuyên biệt \& Tích hợp Agents}
\begin{enumerate}
    \item \textbf{Cài đặt Wazuh Manager trên máy chủ \texttt{wazuh-siem} (10.0.20.10)}: Cài đặt gói \texttt{wazuh-manager=4.14.9-1}. Tắt module \texttt{<vulnerability-detection>} để tối ưu dung lượng đĩa và tài nguyên RAM, bảo đảm máy chủ luôn hoạt động mượt mà và ổn định.
    \item \textbf{Cài đặt và kết nối Wazuh Agent trên Gateway 1, Gateway 2 và K8s Master}:
    \begin{itemize}
        \item Đăng ký Agent tự động qua Ansible role \texttt{wazuh\_agent} kết nối tới IP \texttt{10.0.20.10}.
        \item Cấu hình Agent trên Gateway đọc trực tiếp file nhật ký cảnh báo của Suricata dạng JSON:
        \begin{lstlisting}[language=xml]
<localfile>
  <log_format>json</log_format>
  <location>/var/log/suricata/eve.json</location>
</localfile>
        \end{lstlisting}
        \item Cấu hình Agent trên Master node thu thập nhật ký K8s audit log tại \texttt{/var/log/kubernetes/audit.log}.
    \end{itemize}
\end{enumerate}

\section{Bảo mật hạ tầng Cloud \& Kubernetes Hardening}
\begin{itemize}
    \item \textbf{Bảo vệ thông tin xác thực AWS qua IMDSv2}: Toàn bộ máy chủ EC2 đều cấu hình \texttt{http\_tokens = "required"} và \texttt{http\_put\_response\_hop\_limit = 2}, ngăn chặn hình thức tấn công SSRF đánh cắp IAM Role credentials.
    \item \textbf{Mã hóa đĩa cứng tầng lưu trữ}: 100\% các ổ cứng EBS (Root Volume) đều được mã hóa bằng khóa quản lý AWS KMS (\texttt{encrypted = true}).
    \item \textbf{PodDisruptionBudget}: Áp dụng \texttt{maxUnavailable: 1} cho các dịch vụ cốt lõi (\texttt{frontend}, \texttt{cartservice}, \texttt{checkoutservice}).
    \item \textbf{PodAntiAffinity}: Cấu hình phân tán (\texttt{preferredDuringSchedulingIgnoredDuringExecution}) cho các Pod \texttt{frontend}, đảm bảo các bản sao Pod được ưu tiên xếp đều trên cả 2 Worker node ở 2 Availability Zones khác nhau.
\end{itemize}
